\chapter{Discussion and Policy Implications}
\label{ch:discussion}

The results presented in Chapter \ref{ch:results} provide robust evidence for the adverse incorporation of informal labor within India's textile GVCs. This chapter contextualizes these findings within the broader structuralist literature, derives policy recommendations from the simulation exercise, and outlines the limitations of the current study. Detailed results for all robustness checks and supplementary tests are provided in Appendix \ref{appendixC}.

\section{Defending the Null Result: Monopsony vs. Measurement Error}
\label{sec:null_defense}

A primary finding of this dissertation is the statistical zero pass-through of formal productivity to informal wages ($\beta = -0.014$, $p = 0.610$ in the pooled specification). Critics might argue that this null result is a consequence of \textbf{attenuation bias} or measurement error in the ASUSE informal wage data, which is notoriously difficult to capture accurately in survey-based instruments.

However, I argue that the null result reflects a substantive economic reality for four reasons. First, the same dataset yields a precise and significant estimate for the \textbf{outsourcing volume} effect ($\beta = 9.07 \times 10^{-4}$, $p = 0.043$). If the data were purely noise, we would expect insignificant coefficients across all specifications. The selectivity of significance—strong for outsourcing, null for wages—is itself informative. Second, the null wage result is extraordinarily stable: it survives six distinct robustness checks including outlier exclusion ($\beta = -0.021$, $p = 0.329$) and an alternative median wage measure. Third, as demonstrated in the policy simulation (Section~\ref{sec:simulation_results}), enforcement intensity is fundamentally incapable of shifting informal wages in any direction. Fourth, the pattern is replicated under rigorous identification bounding. Applying Oster-style sensitivity bounds to the primary outsourcing model ($RV = 7.6\%$) demonstrates that an unobserved confounder would need to explain 7.6\% of both the residual variation in outsourcing intensity and in formal productivity to reduce the core elasticity to zero. While this robustness value is modest, it indicates that even a confounder substantially stronger than state enforcement ($E_s$) could not fully overturn the main finding.

I acknowledge, however, that the state-industry level of aggregation may mask within-district heterogeneity. It is possible that wage pass-through occurs locally in tight labor markets (e.g., garment clusters in Tiruppur or Surat) but is averaged out at the state level. District-level microdata, currently suppressed in public ASI-ASUSE releases, would allow a sharper test of this possibility.

\section{Rethinking Supplementary Tests: Inconclusive but Informative}
\label{sec:inconclusive_tests}

Two of the supplementary political economy tests yielded inconclusive results that warrant explicit discussion rather than dismissal.

\textbf{Supplementary Test A: Industrial Density.} The elasticity of informal enterprise counts with respect to formal productivity is $-0.14$ ($p = 0.257$). The negative sign is \textit{inconsistent} with the satellite hypothesis under adverse incorporation, but the lack of significance prevents strong conclusions. Two interpretations are plausible. First, the state-industry sample ($N = 141$) is simply underpowered to detect the relatively small enterprise-count effect. Monte Carlo power analysis suggests that detecting a medium effect size at this sample size achieves only around 55\% power. Second, the enterprise-count measure conflates two opposing forces: formal productivity growth may simultaneously \textit{generate} informal satellites for subcontracting while \textit{reducing} informal self-employment as workers transition. These countervailing channels may cancel at the aggregate level. The persistence test in Section \ref{sec:persistence_test}, which focuses on employment stocks rather than enterprise counts, provides more conclusive evidence for structural embedding precisely because it is less susceptible to this compositional problem.

\textbf{Supplementary Test B: Cost-Shifting via Outsourcing.} The confirmation of the cost-shifting test via orthogonalized mediation ($b = -3.919$, $p = 0.005$; raw mediation: $\beta = -7.226$, $p = 0.009$) provides the ``missing link'' for the adverse incorporation argument. It demonstrates that formal lead firms do not merely coexist with an informal fringe; they actively extract surplus through subcontracting channels. As these firms modernize and expand outsourcing volume beyond their baseline productivity trends, they leverage their buyer power to squeeze the prices paid to informal workshops. This effectively shifts the costs of regulatory compliance and market volatility onto the most precarious segments of the labor force. Notably, the VIF checks on the baseline mediation model confirm acceptable collinearity ($\text{VIF}(\ln A_F)=1.08$, $\text{VIF}(\ln Q_{\text{out}})=1.27$), so the orthogonalization is a rigour exercise rather than a necessity. The sustained negative and significant coefficient after orthogonalization confirms that the ``price squeeze'' mechanism identified in global apparel value chains \autocite{Anner2020} is now directly observable within domestic Indian state-industry networks.

\section{The Persistence Finding as the Strongest Evidence}
\label{sec:persistence_interpretation}

	Among the three research questions, the persistence test (Section \ref{sec:persistence_test}) provides the most unambiguous evidence for structural embedding. The autoregressive coefficient $\beta_1 = 0.8503$ (SE: 0.041, $p < 0.001$) implies that roughly 85\% of a state-industry's informal employment is predetermined by its own history, with formal productivity contributing a statistically and economically negligible increment ($\beta_2 = -0.051$, $p = 0.370$). Crucially, this path dependence persists after controlling for year fixed effects and contemporaneous enforcement intensity, ruling out the possibility that persistence merely reflects stable macro conditions.

The magnitude—0.85—places India's textile informal sector firmly in the ``structural lock-in'' regime identified by \textcite{Portes1989} and is consistent with the \textcite{Venkatesh2006} observation that informal economies reproduce themselves through dense social networks and established subcontracting relationships that are insensitive to market-level productivity shocks. The AR(2) robustness check confirms that this result is not an artifact of omitted higher-order dynamics: the second lag is non-significant ($\beta = 0.253$, $p = 0.185$), validating the AR(1) specification.

I acknowledge the unit root concern candidly. While the LLC panel unit root test could not be executed with sufficient power (only 2 balanced units), the Fisher Combined ADF test yields $p = 0.054$ (borderline at the 10\% level). The block-bootstrapped 95\% BCa confidence interval for the AR(1) coefficient is $[0.734,\; 0.968]$, firmly excluding the unit root ($\beta_1=1$) from below while lying well above the 0.5 structural lock-in threshold. The first-difference specification ($\beta_{FD} = -0.038$, $p = 0.420$) is consistent with a stable structural equilibrium rather than explosive non-stationarity. Future research should prioritize a longer continuous time series for formal unit root rejection.

\section{Endogeneity and Causality}
\label{sec:causality}

The political economy tests implemented in this study identify \textbf{equilibrium correlations} rather than strict causal pathways. A central concern is that state enforcement intensity ($E_s$) is not randomly assigned. States with strong labor institutions—Kerala, West Bengal, Jammu and Kashmir—differ systematically from others in political ideology, historical union density, and industrial composition. These unobserved factors could independently drive higher wages or outsourcing behavior, creating omitted variable bias in the enforcement interaction terms. While instrumental variable strategies would strengthen causal claims, credible instruments for enforcement intensity are unavailable in the ASI-ASUSE public data.

I therefore rely on a triangulation strategy rather than asserting definitive causality. First, the positive outsourcing-productivity elasticity survives six robustness checks, including outlier exclusion, alternative enforcement measures, non-linear enforcement specifications, and multi-year replication. Second, as shown via the Oster-style bounding exercise, an unobserved confounder would need to explain 7.6\% of both the residual variation in outsourcing \textit{and} in formal productivity to fully overturn the main elasticity result. While this robustness value is modest by the standards of experimental economics, it implies that a confound substantially larger in explanatory power than state enforcement is required. The enforcement heterogeneity finding reflects a genuine structural pattern, but in the absence of a clean natural experiment, I do not claim that formal productivity strictly \textit{causes} wage stagnation in a vacuum; rather, the data \textit{suggests} that both the outsourcing expansion and the wage null are outcomes of a unified GVC strategy in which formal lead firms conditionally exploit buyer power to capture productivity rents.

\section{Policy Implications}
\label{sec:policy_implications}

The empirical findings and simulation results lead to three critical policy recommendations for labor governance in India's dual economy.

\textbf{Stability over a 15-Year Horizon:} By extending the panel from 2010 to 2024, I show that these patterns are not a temporary transitional friction but a stable structural equilibrium. The persistence coefficient ($\beta=0.850$) remains remarkably high even through massive economic shocks, including the 2016 demonetization and the COVID-19 pandemic. This suggests that the informal-formal link is robustly embedded in Indian manufacturing's structural DNA.
\textbf{1. The Primacy of the Wage Floor.} The policy simulation (Section \ref{sec:simulation_results}) demonstrates that conventional enforcement shocks are fundamentally incapable of raising informal wages. Stricter enforcement merely scales the informal sector without improving the terms of incorporation. The only intervention that successfully broke the monopsony trap in the simulation was an exogenous \textbf{direct wage floor}. This suggests that policy should shift from procedural compliance (inspecting factory floors) to substantive distributional mandates (statutory minimum wages for job workers) that bypass the formal buyer's piece-rate setting power.

\textbf{2. The Perverse Incentives of Enforcement.} The finding that high-enforcement states exhibit 4.4 times greater outsourcing elasticity illustrates a "regulatory arbitrage" trap. When labor laws are enforced strictly \textit{within} the formal factory gates but not \textit{across} the value chain, enforcement perversely incentivizes displacement to the unregulated informal sector. This "scaling without improvement" effect suggests that traditional inspectorate-based enforcement may be counterproductive if it remains geographically and sectorally narrow.

\textbf{3. Beyond the Factory Gate: Value Chain Accountability.} If the "price squeeze" is the primary mechanism for wage suppression, then regulation must move from the production site to the \textbf{contractual relationship}. Policy frameworks like Germany's Supply Chain Act (LKSG) offer a theoretical blueprint: formal lead firms must be held legally accountable for labor standards and fair pricing across their subcontracting networks. By internalizing the social reproduction costs of informal labor, such regulations could neutralize the incentive for adverse incorporation.

\textbf{4. Feasibility Constraints in Expanding Minimum Wages.} While a direct wage floor is the only mathematically effective intervention in the simulation, implementing such a policy mechanism for decentralized, home-based workers faces severe administrative constraints in India. The current enforcement inspectorate ratio ($\approx 0.04$ inspectors per registered factory in major states like Gujarat and Maharashtra) is already vastly overwhelmed by the formal sector alone. Attempting to enforce statutory piece-rates across millions of unregistered household enterprises would require an unprecedented expansion of state capacity and fiscal expenditure. Furthermore, the legal classification of "job worker" vs "independent contractor" complicates enforcement, creating loopholes for lead firms to deny employment status. Policy must therefore focus on structural upstream interventions—targeting the centralized formal buyers who control the payment terms—rather than attempting to police the fragmented informal downstream network.

\section{Limitations}
\label{sec:limitations_discussion}

Five limitations warrant emphasis in addition to those discussed in Section \ref{sec:limitations} of Chapter \ref{ch:results}.

\textbf{Sample size and power.} The cross-sectional regression sample of $N = 47$ state-industry cells is small by the standards of applied microeconometrics. While the outsourcing result ($p = 0.043$) survives six robustness checks, the interaction terms and heterogeneity analyses are underpowered. The industrial density and cost-shifting tests remain inconclusive, and negative findings in these tests should be interpreted as low-power null results rather than evidence against adverse incorporation.

\textbf{Geographic suppression.} Public ASI data suppresses district codes, forcing aggregation to the state level. This masks within-state variation in labor market tightness—arguably the most theoretically relevant source of identification for wage pass-through. District-level data, accessible only through restricted-access ASI microfiles, would allow a sharper test of local labor market effects.

\textbf{Panel length.} Three years of panel data is sufficient for the persistence test but marginal for formal unit root diagnostics. A five-to-seven year panel would permit the LLC test, provide more precise autoregressive estimates, and allow the identification of structural breaks coinciding with policy reforms (e.g., the 2019 Wage Code amendments).

\textbf{Task composition.} Aggregate outsourcing costs (Block F) do not distinguish which tasks are subcontracted. If high-productivity formal firms systematically outsource lower-skill tasks, the observed wage null could reflect compositional downgrading rather than pure monopsony rent-capture. Disentangling these mechanisms requires task-level data currently unavailable in public surveys.

\textbf{Firm-level matching.} The ASI and ASUSE cannot be matched at the firm level due to anonymization. A matched formal-informal dataset—linking lead firms to their specific subcontractors—would allow a direct test of the price-squeeze mechanism and would substantially sharpen inference on the monopsony hypothesis.

\section{Future Research}
\label{sec:future}

Three research extensions follow naturally from these findings. First, district-level analysis using restricted-access ASI microfiles would allow wage pass-through tests in local labor markets where informality is concentrated, addressing the geographic suppression problem. Second, longitudinal tracking of informal enterprises—ideally through a panel component added to the ASUSE—would allow direct tests of whether informal units graduate to formal status or remain trapped as precariously paid suppliers over time, providing a more definitive resolution of the transitional versus structural debate. Third, integrating buyer-supplier transaction data from GST filings (which record inter-firm payments at the firm level) with ASUSE informal enterprise data would permit a direct test of the price-squeeze mechanism, sidestepping the collinearity problem that limited Supplementary Test B in the current analysis.
